\chapter{Sorting, Selection and Sorting Networks}\label{ch:sorting}

\begin{center}\itshape
A network of comparators that sorts every sequence of $0$s and $1$s sorts every sequence of numbers.\\
\upshape(the zero--one principle, in our words)
\end{center}

In the spring of 1945, while the first stored-program computer existed only on paper, John von Neumann wrote one of the first programs
ever written for such a machine. He did not choose a numerical problem, although numerical problems were what the machine was being
built for. He chose \emph{sorting}: twenty-three handwritten pages, for the EDVAC, implementing the merging of two sorted lists, the heart of
merge sort. He wanted to show that a general-purpose computer could beat the special punched-card sorting machines of the day at their own
game. A quarter of a century later Donald Knuth found the manuscript and analysed it, bug included.

Fourteen years later a young Englishman named Tony Hoare was a visiting student at Moscow State University. He had been offered a job
on machine translation, and to translate a sentence one looks up its words in a dictionary stored, in alphabetical order, on magnetic
tape: so the words of the sentence must first be sorted. He thought first of the obvious method, found it slow, and then came up with
\emph{quicksort}: choose one element, the pivot, put the smaller ones to its left and the larger ones to its right, and sort the two sides in
the same way. Published in 1961, it is still, in one form or another, the method inside the sorting routines of most programming
languages.

This chapter is about putting things in order, and about the cheaper task of finding one element of the order, such as the median. We
compute merge sort's worst case exactly, and prove that quicksort, whose average was found in Tome~3, is very unlikely to be much
worse than its average. We find the median in linear time, deterministically by a beautiful recursion and randomly by sampling, with
fewer comparisons than any deterministic method can achieve. We sort without comparing at all, by digits. And we build sorting machines
whose comparisons are fixed in advance and run in parallel: the \emph{sorting networks}, among them Batcher's network on the cover of this
book, whose correctness rests on a principle with a one-line proof: to check that a network sorts everything, it is enough to check $0$s and $1$s.

\begin{think}
Twenty-five horses, a racetrack with five lanes and no stopwatch: each race ranks five horses. How many races do you need to find the
three fastest horses? Seven suffice: can you find them? And can you show that five races cannot even find the fastest horse?
\end{think}

\section{Merge sort, exactly}\label{sec:mergesort}

Merge sort splits the list into halves of sizes $\lceil n/2\rceil$ and $\lfloor n/2\rfloor$, sorts them recursively, and merges them. Proposition~1.8 showed
that it never needs more than $n\lceil\log_2n\rceil$ comparisons. Now we find its worst case exactly.

\begin{lemma}\label{lem:mergecost}
Merging sorted lists of lengths $p,q\ge1$ by repeatedly comparing the two front elements makes $p+q-r$ comparisons, where $r$ is the number of
elements left in one list when the other is exhausted. It makes $p+q-1$ comparisons if and only if the two largest elements lie in
different lists.
\end{lemma}
\begin{proof}
Each comparison outputs one element; when a list is exhausted the remaining $r\ge1$ elements are copied without comparison. So there
are $p+q-r$ comparisons. Now $r=1$ exactly when the last element output by a comparison is the second largest and the largest remains alone:
that is, when the largest element is in one list and the second largest in the other. (If both are in the same list, that list still holds
both when the other is exhausted, so $r\ge2$.)
\end{proof}

\begin{theorem}\label{thm:mergeexact}
The worst-case number of comparisons of merge sort on $n\ge1$ elements is
\[
W(n)=\sum_{k=1}^n\lceil\log_2k\rceil=n\lceil\log_2n\rceil-2^{\lceil\log_2n\rceil}+1.
\]
\end{theorem}
\begin{proof}
\emph{The recurrence.} We show $W(1)=0$ and $W(n)=W(\lceil n/2\rceil)+W(\lfloor n/2\rfloor)+n-1$ for $n\ge2$. The comparisons made inside each half depend only
on the relative order of the elements of that half, and the merge makes at most $n-1$ comparisons, which gives $\le$. For $\ge$, take inputs of
sizes $\lceil n/2\rceil$ and $\lfloor n/2\rfloor$ achieving the worst cases of the halves, and choose the actual values so that each half keeps its relative order and
the largest element of the first half is the overall maximum while the largest of the second half is the second largest. Then the halves
cost $W(\lceil n/2\rceil)+W(\lfloor n/2\rfloor)$, and the merge costs $n-1$ by Lemma~\ref{lem:mergecost}.

\emph{The sum satisfies the recurrence.} Let $g(n)=\sum_{k=1}^n\lceil\log_2k\rceil$ and $h(n)=g(\lceil n/2\rceil)+g(\lfloor n/2\rfloor)+n-1$, with $g(0)=0$; then $h(1)=0=g(1)$. We show
$h(n)-h(n-1)=\lceil\log_2n\rceil=g(n)-g(n-1)$ for $n\ge2$, and then $h=g$, hence $W=g$ by induction. If $n=2m$, only $\lfloor n/2\rfloor$ changes, from $m-1$ to $m$,
so $h(n)-h(n-1)=\lceil\log_2m\rceil+1=\lceil\log_2(2m)\rceil$. If $n=2m+1$ with $m\ge1$, only $\lceil n/2\rceil$ changes, from $m$ to $m+1$, so $h(n)-h(n-1)=\lceil\log_2(m+1)\rceil+1=\lceil\log_2(2m+2)\rceil$, and
this equals $\lceil\log_2(2m+1)\rceil$ because $\lceil\log_2x\rceil$ only jumps just after a power of $2$, and the odd number $2m+1>1$ is not one.

\emph{The closed form.} With $t=\lceil\log_2n\rceil$, the integers $k\le n$ with $\lceil\log_2k\rceil=i$ are those in $(2^{i-1},2^i]$, so
\[
g(n)=\sum_{i=1}^{t-1}i\,2^{i-1}+t\,(n-2^{t-1})=\bigl((t-2)2^{t-1}+1\bigr)+tn-t2^{t-1}=tn-2^t+1,
\]
using $\sum_{i=1}^{t-1}i2^{i-1}=(t-2)2^{t-1}+1$, an easy induction.
\end{proof}

Compare with the information bound $\lceil\log_2n!\rceil$ of Theorem~1.15: since $\log_2n!=\sum_{k\le n}\log_2k$, merge sort pays, for each $k$, the
round-up of $\log_2k$ instead of $\log_2k$. The loss is less than $n$ comparisons in all, and in relative terms it tends to $0$.

\section{Quicksort}\label{sec:quicksort}

\emph{Randomised quicksort} chooses a pivot uniformly at random in the current sublist, compares it with every other element of the
sublist, splits them into the smaller and the larger ones, and recurses on both sides with fresh random choices (Figure~\ref{fig:quicksort}). Its
worst case is bad: if every pivot happens to be the smallest element, it makes $\binom n2$ comparisons. But Tome~3, Exercise~2.3, showed that two
elements of ranks $i<j$ are compared with probability exactly $2/(j-i+1)$, so that the \emph{expected} number of comparisons is
\[
\mathbb E[C_n]=2(n+1)H_n-4n=2n\ln n+O(n)\approx1.39\,n\log_2n,
\]
where $H_n=1+\frac12+\dots+\frac1n$. So on average quicksort uses about $39\%$ more comparisons than merge sort; it wins in practice because each
comparison is cheap (the pivot sits in a register, the elements are scanned in order) and it sorts in place.

\begin{figure}[ht]
\centering
\fig{ch02-quicksort}
\caption{Quicksort on nine elements, with the first element of each sublist as pivot (in orange). Each element is compared with the
pivot of every sublist it belongs to, except its own: the total number of comparisons is the sum, over the elements, of their depth in this
tree. Here it is $17$, below $\lceil\log_29!\rceil=19$: on one input an algorithm can be lucky, and the bound of Theorem~1.15
concerns the worst case and the average.}
\label{fig:quicksort}
\end{figure}

An average says nothing about the chance of a disaster. The tree of Figure~\ref{fig:quicksort} suggests how to control it: the number of comparisons
is the sum of the depths of the elements, and an element sinks fast as soon as some pivots on its path are reasonably central. We need a way to
compare a process whose steps succeed with probability \emph{at least} $\frac12$, but not independently, with independent coin flips.

\begin{lemma}\label{lem:dominate}
Let $Y_1,\dots,Y_L$ be random variables with values in $\{0,1\}$ such that, for each $d$, conditionally on any values of $Y_1,\dots,Y_{d-1}$ (and of
any other information revealed before step $d$), $P(Y_d=1)\ge\frac12$. Then, for every $G$, $P(Y_1+\dots+Y_L<G)\le P(B<G)$ where $B$ is binomial with
parameters $L$ and $\frac12$.
\end{lemma}
\begin{proof}
We build the process from independent uniform variables $U_1,\dots,U_L$ on $[0,1]$. Given the past, let $p_d\ge\frac12$ be the conditional
probability that $Y_d=1$, and generate step $d$ by setting $Y_d=1$ if $U_d\le p_d$, the rest of the step being drawn from its conditional law given
$Y_d$ (with fresh randomness). This produces a process with the same law. Now $B_d=\mathbf 1[U_d\le\frac12]$ are independent Bernoulli$(\frac12)$ and
$B_d\le Y_d$ for every $d$. So $Y_1+\dots+Y_L<G$ implies $B_1+\dots+B_L<G$.
\end{proof}

\begin{theorem}\label{thm:qsconc}
For $n\ge2$, with probability at least $1-n^{-2}$, randomised quicksort on $n$ elements makes at most $n\lceil30\ln n\rceil$ comparisons, and its recursion
has depth at most $\lceil30\ln n\rceil+1$.
\end{theorem}
\begin{proof}
Fix an element $x$, and let $s_0=n,s_1,s_2,\dots$ be the sizes of the successive sublists containing $x$, as long as $x$ has not been chosen as a
pivot and the sublist has size at least $2$. At step $d$ the element $x$ is compared once, with the pivot. Call step $d$ \emph{good} if both parts
produced by the pivot have size at most $\frac34s_d$; a pivot of rank $r$ in the sublist leaves parts of sizes $r-1$ and $s_d-r$, so the step is good
when $\frac14s_d\le r\le\frac34s_d+1$, which happens for at least $s_d/2$ of the $s_d$ equally likely ranks. Whatever happened before, a step is good with
probability at least $\frac12$; once $x$'s journey is over, declare all further steps good.

After $G$ good steps the sublist of $x$ has size at most $(\frac34)^Gn$, which is less than $2$ once $G\ge\log_{4/3}n$. Let $G=\lceil\log_{4/3}n\rceil\le\frac{\ln n}{\ln(4/3)}+1\le5\ln n$
(as $1/\ln\frac43<3.48$ and $1\le1.5\ln n$ for $n\ge2$), and $L=\lceil30\ln n\rceil$. If $x$ is compared more than $L$ times, then fewer than $G$ of the first $L$ steps are
good. By Lemma~\ref{lem:dominate} and the lower Chernoff bound (Tome~3, Theorem~4.6) applied to $B$, with mean $\mu=L/2\ge15\ln n$ and $\delta=\frac23$, for which $(1-\delta)\mu\ge5\ln n\ge G$:
\[
P(\text{$x$ is compared more than $L$ times})\le P\bigl(B\le(1-\delta)\mu\bigr)\le e^{-\delta^2\mu/2}\le e^{-\frac49\cdot\frac{15}2\ln n}=n^{-10/3}.
\]
By the union bound over the $n$ elements, with probability at least $1-n^{-7/3}\ge1-n^{-2}$ no element is compared more than $L$ times as a non-pivot.
Every comparison involves one non-pivot, so there are at most $nL$ comparisons. And an element at depth $D$ of the recursion was a non-pivot in
$D-1$ sublists of size at least $2$ before, so the depth is at most $L+1$.
\end{proof}

The constant $30$ is lavish, and the truth is much better: $C_n/(2n\ln n)\to1$ in probability, and the fluctuations of $C_n$ around its mean are of
order $n$ (R\'egnier, R\"osler, 1989--1991). The notebook draws the distribution. The theorem is enough to say that quicksort's bad cases are not just
rare but astronomically rare: for $n=10^6$ the probability of exceeding $30n\ln n$ comparisons is below $10^{-12}$.

\begin{think}
A deterministic quicksort that always takes the \emph{first} element as pivot has worst case $\binom n2$, and the worst input is a sorted list,
which is not rare in practice. Randomising the pivot makes every input behave like a random input. Is there still a ``bad input'' for
randomised quicksort? What exactly is random in Theorem~\ref{thm:qsconc}: the input or the algorithm?
\end{think}

\section{Selection}\label{sec:selection}

To find the $k$-th smallest of $n$ elements (the median when $k=\lceil n/2\rceil$) one can sort, in $O(n\log n)$. One can do better, because most of the
order is irrelevant.

\subsection*{Quickselect}

Hoare's \emph{quickselect} partitions around a random pivot like quicksort, but recurses only into the side that contains the $k$-th element (or
stops if the pivot is it).

\begin{proposition}\label{prop:quickselect}
For every $k$, quickselect makes at most $8n$ comparisons on average.
\end{proposition}
\begin{proof}
Say the algorithm is in \emph{phase} $j\ge0$ while the current sublist has size $s$ with $(\frac34)^{j+1}n<s\le(\frac34)^jn$. Each round costs $s-1<(\frac34)^jn$
comparisons. As in the proof of Theorem~\ref{thm:qsconc}, with probability at least $\frac12$, whatever happened before, the pivot leaves two parts of size at most
$\frac34s$; then the next sublist, whichever it is, has size at most $\frac34s\le(\frac34)^{j+1}n$, and phase $j$ is over. So the number $R_j$ of rounds spent in phase $j$
satisfies $P(R_j\ge t)\le2^{-(t-1)}$, and $\mathbb E[R_j]=\sum_{t\ge1}P(R_j\ge t)\le2$. The expected number of comparisons is at most
$\sum_j(\frac34)^jn\,\mathbb E[R_j]\le2n\sum_{j\ge0}(\frac34)^j=8n$.
\end{proof}

The exact average is known (Knuth, 1971): for the median it is about $3.39n$, for the minimum about $2n$ (Exercise~\ref{ex:qselectmin}).

\subsection*{The median of medians}

In 1973 five authors, Manuel Blum, Robert Floyd, Vaughan Pratt, Ronald Rivest and Robert Tarjan (four of them future Turing Award winners),
found a deterministic linear-time method. The idea is to choose a pivot that is guaranteed to be central, by a recursive call on a small sample.

\begin{theorem}[Blum, Floyd, Pratt, Rivest and Tarjan, 1973]\label{thm:bfprt}
The following algorithm $\mathrm{Select}(S,k)$ finds the $k$-th smallest element of a set $S$ of $n$ elements in time $O(n)$:
if $n<140$, sort $S$ and answer; otherwise split $S$ into $\lceil n/5\rceil$ groups of at most $5$, find the median of each group, compute recursively the median $m$ of
these medians, split $S$ into $S_<$, $\{m\}$, $S_>$ by comparing with $m$, and recurse into the part containing the answer.
\end{theorem}
\begin{proof}
\emph{Correctness} is clear: $m\in S$, and the answer is in the part that has the right rank range.

\emph{$m$ is central.} At least half of the $g=\lceil n/5\rceil$ group medians, that is at least $\lceil g/2\rceil$ of them, are $\le m$. Each such median, except possibly that of the
incomplete group and that of $m$'s own group, comes with at least $3$ elements of its full group that are $\le$ it, hence $\le m$. So at least
$3(\lceil g/2\rceil-2)\ge\frac{3n}{10}-6$ elements are $\le m$, and therefore $|S_>|\le\frac{7n}{10}+6$; symmetrically $|S_<|\le\frac{7n}{10}+6$.

\emph{Time.} Finding the group medians and splitting cost at most $an$ for a constant $a$. Let $T(n)$ be the worst-case time; then $T(n)\le C_0$ for $n<140$ and
$T(n)\le T(\lceil n/5\rceil)+T(\lfloor7n/10\rfloor+6)+an$ for $n\ge140$, with $T$ taken nondecreasing (replace $T(n)$ by $\max_{m\le n}T(m)$). Let $c=\max(C_0,20a)$. We show $T(n)\le cn$ by
induction. For $n<140$ it holds since $c\ge C_0$. For $n\ge140$, both recursive sizes are smaller than $n$, and
\[
T(n)\le c\Bigl(\frac n5+1\Bigr)+c\Bigl(\frac{7n}{10}+6\Bigr)+an\le cn-\Bigl(\frac{cn}{10}-7c-\frac{cn}{20}\Bigr)=cn-c\Bigl(\frac n{20}-7\Bigr)\le cn.
\]
\end{proof}

This is Proposition~1.9(i) in action: the two recursive calls have total size about $(\frac15+\frac7{10})n=\frac9{10}n<n$, so the work shrinks geometrically down the recursion.
Groups of $3$ would give $\frac13+\frac23=1$, and only $O(n\log n)$ (Exercise~\ref{ex:groups}).

\subsection*{Selection by sampling}

The median of medians is a theoretical jewel, but its constant factor is large. Bent and John proved in 1985 that every
deterministic algorithm needs at least $2n-o(n)$ comparisons to find the median. Randomisation beats this: a random sample tells us where the median
roughly is, and only the few elements near it need careful treatment. This idea goes back to Floyd and Rivest (1975); we prove the clean version of
Motwani and Raghavan.

\begin{theorem}\label{thm:lazyselect}
Let $S=\{s_1<s_2<\dots<s_n\}$, $n\ge4$, and $1\le k\le n$. One \emph{trial} of the following algorithm fails with probability at most $8n^{-1/4}$, and if it does not fail it
returns $s_k$.
\begin{enumerate}
\item Draw $r=\lceil n^{3/4}\rceil$ elements of $S$ independently and uniformly, with replacement, and sort them: $R_{(1)}\le\dots\le R_{(r)}$.
\item Let $x=kr/n$, $\ell=\lfloor x-\sqrt n\rfloor$, $h=\lceil x+\sqrt n\rceil$. Let $a=R_{(\ell)}$ if $\ell\ge1$ and $a=-\infty$ otherwise; $b=R_{(h)}$ if $h\le r$ and $b=+\infty$ otherwise.
\item Compare every element of $S$ with $a$ and $b$; let $P=\{y\in S:a\le y\le b\}$ and $t=|\{y\in S:y<a\}|$.
\item If $t<k\le t+|P|$ and $|P|\le4n^{3/4}+2$, sort $P$ and return its $(k-t)$-th element; otherwise the trial \emph{fails}.
\end{enumerate}
Repeating until success, the expected number of comparisons is $n+\min(k,n-k)+O(n^{3/4}\log n)$.
\end{theorem}
\begin{proof}
If the trial does not fail, $s_k\in P$ and it is the $(k-t)$-th element of $P$, so the answer is right. Put $s_j=-\infty$ for $j\le0$ and $s_j=+\infty$ for $j>n$, and let
$k'=k-\lceil2n^{3/4}\rceil$ and $k''=k+\lceil2n^{3/4}\rceil$. The trial can only fail if one of four events happens:
\[
\begin{aligned}
&E_1:\ s_k<a,\qquad E_2:\ s_k>b,\\
&E_3:\ k'\ge1\text{ and }a\le s_{k'},\qquad E_4:\ k''\le n\text{ and }b\ge s_{k''}.
\end{aligned}
\]
Indeed, if none happens then $a\le s_k\le b$, so $t<k\le t+|P|$; and every element of $P$ is some $s_j$ with $j>k'$ (if $k'\ge1$ because $a>s_{k'}$, and trivially otherwise) and
$j<k''$, so $|P|\le k''-k'-1\le2\lceil2n^{3/4}\rceil-1\le4n^{3/4}+1$.

Each event says that a binomial count of sample elements is far from its mean. We use Chebyshev's inequality (Tome~3, Proposition~4.1) in the form: a binomial
variable $Z$ with $r$ trials has $\operatorname{Var}Z\le r/4$, so $P(|Z-\mathbb EZ|\ge\frac{\sqrt n}2)\le\frac{r/4}{n/4}=\frac rn\le\frac{2n^{3/4}}n=2n^{-1/4}$.

$E_1$: then $\ell\ge1$ and fewer than $\ell$ sample elements are $\le s_k$. Their number $Z$ is binomial with mean $kr/n=x$, and $Z<\ell\le x-\sqrt n$.

$E_2$: then $h\le r$ and at least $h$ sample elements are $<s_k$. Their number $Z$ has mean $(k-1)r/n\le x$, and $Z\ge h\ge x+\sqrt n$.

$E_3$: then $k'\ge1$. This forces $\ell\ge1$: if $\ell<1$ then $x<1+\sqrt n$, so $k<n(1+\sqrt n)/r\le n^{1/4}+n^{3/4}\le\lceil2n^{3/4}\rceil$ and $k'\le0$. So $a$ is a sample element, and at
least $\ell$ sample elements are $\le s_{k'}$; their number $Z$ has mean $k'r/n\le x-2n^{3/4}r/n\le x-2\sqrt n$, and $Z\ge\ell>x-\sqrt n-1$, so $Z-\mathbb EZ>\sqrt n-1\ge\frac{\sqrt n}2$
(as $n\ge4$).

$E_4$: symmetric to $E_3$, exchanging the roles of small and large elements.

Each event has probability at most $2n^{-1/4}$, so a trial fails with probability at most $8n^{-1/4}$.

\emph{Comparisons.} Sorting the sample and $P$ costs $O(n^{3/4}\log n)$. In step 3, if $k\le n/2$ compare each $y$ first with $b$, and with $a$ only when $y<b$; this
costs $n+|\{y:y<b\}|$, which on success is at most $n+k+|P|$. If $k>n/2$, start with $a$, for $n+(n-k)+|P|+1$. A failed trial costs at most $2n+O(n^{3/4}\log n)$, and the
number of trials is geometric with expectation at most $1/(1-8n^{-1/4})=1+O(n^{-1/4})$. So the expected total is $n+\min(k,n-k)+O(n^{3/4}\log n)+O(n^{-1/4})\cdot O(n)$, as claimed.
\end{proof}

For the median that is $1.5n+o(n)$ comparisons, well below the deterministic $2n$. (For $n$ below about $4100$ the bound $8n^{-1/4}$ exceeds $1$ and says nothing;
in practice the method works for much smaller $n$, as the notebook shows, and production versions use a larger sample.)

\section{Sorting without comparing}\label{sec:radix}

The lower bound $\log_2n!$ is about comparisons. If the keys are small integers we can use them as addresses, and the bound disappears.

\begin{proposition}[Counting sort]\label{prop:counting}
Records with keys in $\{0,1,\dots,K-1\}$ can be sorted \emph{stably} (records with equal keys keep their input order) in time $O(n+K)$.
\end{proposition}
\begin{proof}
Count the records of each key, $c[0..K-1]$, in one pass. Compute the starting positions $p[j]=c[0]+\dots+c[j-1]$. Scan the input in order and put each record with key $j$
at position $p[j]$, then increase $p[j]$. Records of key $j$ fill positions $c[0]+\dots+c[j-1]$ to $c[0]+\dots+c[j]-1$, in input order: the output is sorted and stable. Each pass costs
$O(n)$ or $O(K)$.
\end{proof}

\begin{theorem}[Radix sort]\label{thm:radix}
Let the keys be integers in $[0,B^d)$, written with $d$ digits in base $B$. Sorting stably by the last digit, then by the next to last, and so on up to the first digit,
sorts the records. With counting sort for each pass, the time is $O(d(n+B))$.
\end{theorem}
\begin{proof}
We show by induction on $i$ that after $i$ passes the records are sorted by the number formed by their last $i$ digits. For $i=1$ this is the first pass. After pass $i+1$
the records are sorted by digit $i+1$ from the right; two records with the same such digit keep, by stability, the order of pass $i$, which sorted them by their last $i$ digits.
So they are sorted by the last $i+1$ digits, read as a number (lexicographically, most significant first). After $d$ passes they are sorted by the keys.
\end{proof}

With $B=n$, $n$ integers below $n^c$ are sorted in $O(cn)$ time: linear for keys of size polynomial in $n$, which no comparison sort can do (Exercise~\ref{ex:radixbase}). The
punched-card sorting machines of the first half of the twentieth century worked exactly this way, one column per pass, starting with the last digit, which surprises
everyone who first sees it: the most significant digit is sorted last.

\section{Sorting networks}\label{sec:networks}

All the algorithms so far decide which elements to compare next according to the previous answers. A \emph{sorting network} fixes its comparisons in advance. It has $n$
horizontal \emph{wires} carrying values from left to right, and \emph{comparators} $[i:j]$, $i<j$, each taking the values on wires $i$ and $j$ and putting the smaller on wire $i$ and
the larger on wire $j$ (Figure~\ref{fig:bitonic}). Comparators touching disjoint wires can work at the same time, so a network is a parallel algorithm; its \emph{depth} is the
largest number of comparators on a path from input to output, its \emph{size} the number of comparators. A network is also a very simple circuit, which is why the same
designs sort data in graphics cards and route packets in switching networks.

\begin{figure}[ht]
\centering
\fig{ch02-bitonic}
\caption{Batcher's bitonic sorting network on $8$ wires: $24$ comparators in $6$ parallel rounds (a vertical segment between two dots is a comparator, the smaller
value leaving on the upper wire). It sorts the pairs, then merges pairs into sorted fours, then fours into eights, each merge being a ``flip'' round followed by
half-cleaners. The input on the left comes out as $0,1,\dots,7$. The cover shows the same network on $16$ wires.}
\label{fig:bitonic}
\end{figure}

How do we know that a network sorts? There are $n!$ orderings to try. The \emph{zero--one principle}, as Knuth calls it, reduces them to the $2^n$ sequences of $0$s and $1$s,
and, more importantly, gives a way to \emph{think} about networks.

\begin{lemma}\label{lem:monotone}
Let $f$ be a nondecreasing function. If a network transforms $(x_1,\dots,x_n)$ into $(y_1,\dots,y_n)$, it transforms $(f(x_1),\dots,f(x_n))$ into $(f(y_1),\dots,f(y_n))$.
\end{lemma}
\begin{proof}
For a single comparator, $\min(f(u),f(v))=f(\min(u,v))$ and $\max(f(u),f(v))=f(\max(u,v))$ because $f$ preserves order. For a network, apply this comparator by comparator.
\end{proof}

\begin{theorem}[Zero--one principle]\label{thm:zeroone}
If a network sorts every input in $\{0,1\}^n$, it sorts every input of $n$ numbers.
\end{theorem}
\begin{proof}
Suppose the network does not sort $x$: its output $y$ has $y_i>y_{i+1}$ for some $i$. Let $f(t)=1$ if $t\ge y_i$ and $f(t)=0$ otherwise; $f$ is nondecreasing. By the lemma, the network
transforms the $0$--$1$ input $f(x)$ into $f(y)$, which has $f(y_i)=1$ and $f(y_{i+1})=0$: not sorted.
\end{proof}

\subsection*{Batcher's bitonic network}

Kenneth Batcher, at the Goodyear Aerospace Corporation, published two families of sorting networks in 1968. Both merge, in the spirit of merge sort, and both are
proved by the zero--one principle. Call a $0$--$1$ sequence $z_1,\dots,z_N$ \emph{bitonic} if, read around a circle, it changes between $0$ and $1$ at most twice: its $1$s, if any, form a
cyclic interval, and so do its $0$s. Examples: $0011100$, $1100011$, $0000$, $0111$. Reversing a bitonic sequence gives a bitonic sequence. A \emph{half-cleaner} on $N=2m$ wires
is the round of comparators $[i:i+m]$, $1\le i\le m$.

\begin{lemma}\label{lem:halfclean}
If the input of a half-cleaner is bitonic, then both halves of the output are bitonic, every value in the upper half is $\le$ every value in the lower half, and one of the two
halves is \emph{clean} (constant).
\end{lemma}
\begin{proof}
Work modulo $2m$. Suppose first that the $1$s of $z$ form a cyclic interval $I$ of length $|I|\le m$. Two positions $i$ and $i+m$ are at cyclic distance $m$, and an interval of
length at most $m$ cannot contain both; so no comparator sees two $1$s, and the upper half, the minima, is all $0$: clean. The lower half has a $1$ at position $i$ exactly when $i$ or
$i+m$ lies in $I$; the map $i\mapsto i\bmod m$ is injective on $I$ (again because $|I|\le m$) and sends the cyclic interval $I$ of $\mathbb Z/2m$ onto a cyclic interval of $\mathbb Z/m$. So the lower
half is bitonic. Upper all $0$ is $\le$ everything. If $|I|>m$, the $0$s form a cyclic interval of length less than $m$, and the same argument with $0$ and $1$, min and max, upper
and lower, exchanged shows that the lower half is all $1$ and the upper half is bitonic.
\end{proof}

A \emph{bitonic sorter} $\mathrm{BS}(N)$, for $N$ a power of $2$, is a half-cleaner followed by two copies of $\mathrm{BS}(N/2)$, on the upper and the lower halves; $\mathrm{BS}(1)$ is empty.
By induction, with the lemma, it sorts every bitonic $0$--$1$ input: after the half-cleaner both halves are bitonic, get sorted, and the upper half is below the lower one. Its depth
is $\log_2N$.

\begin{theorem}[Batcher, 1968]\label{thm:bitonic}
For $n=2^p$, the network $\mathrm{Sort}(n)$ defined by: $\mathrm{Sort}(1)$ empty; $\mathrm{Sort}(n)$ is two copies of $\mathrm{Sort}(n/2)$ on the two halves, followed by the \emph{flip}
round of comparators $[i:n+1-i]$, $1\le i\le n/2$, followed by a copy of $\mathrm{BS}(n/2)$ on each half, sorts every input. It has depth $\frac12p(p+1)$ and size $\frac n4p(p+1)$.
\end{theorem}
\begin{proof}
By the zero--one principle it suffices to sort $0$--$1$ inputs, and by induction the two halves are sorted when the flip round starts: $x=0^a1^{m-a}$ on wires $1..m$ and
$y=0^c1^{m-c}$ on wires $m+1..2m$, with $m=n/2$. Read the wires in the order $1,2,\dots,m,2m,2m-1,\dots,m+1$: the sequence $w=x\cdot\mathrm{reverse}(y)=0^a1^{2m-a-c}0^c$ is bitonic.
In this order, the comparator $[i:n+1-i]$ compares $w_i$ with $w_{m+i}$, putting the minimum on the upper wire: the flip round is a half-cleaner for $w$. By Lemma~\ref{lem:halfclean}
the upper half (wires $1..m$) and the lower half (wires $2m,\dots,m+1$ in this order, hence also in the natural order, since reversal preserves bitonicity) are bitonic, and
the upper half is below the lower half. The two bitonic sorters finish the job.

Depth: $D(n)=D(n/2)+1+(p-1)=D(n/2)+p$, so $D(2^p)=1+2+\dots+p$. Size: each round has $n/2$ comparators, and there are $D(n)$ rounds.
\end{proof}

For $n=16$, $10$ rounds and $80$ comparators: the network on the cover. Batcher's other network merges differently and saves comparators.

\begin{theorem}[Batcher's odd--even merge]\label{thm:oddeven}
Let $m$ be a power of $2$ and let $x_1\le\dots\le x_m$ and $y_1\le\dots\le y_m$. Merge recursively the odd-indexed elements $x_1,x_3,\dots$ with $y_1,y_3,\dots$ into $v_1\le\dots\le v_m$, and the
even-indexed ones into $w_1\le\dots\le w_m$ (for $m=1$, merging is one comparator). Then the sequence
\[
\begin{aligned}
&v_1,\ \min(w_1,v_2),\ \max(w_1,v_2),\ \min(w_2,v_3),\ \max(w_2,v_3),\\&\dots,\ \min(w_{m-1},v_m),\ \max(w_{m-1},v_m),\ w_m
\end{aligned}
\]
is the sorted merge of $x$ and $y$.
\end{theorem}
\begin{proof}
The whole procedure is a comparator network, so by the zero--one principle we may assume $0$--$1$ inputs. Let $x$ have $a$ zeros and $y$ have $b$ zeros. The odd-indexed elements of $x$
contain $\lceil a/2\rceil$ zeros and the even-indexed ones $\lfloor a/2\rfloor$, and similarly for $y$. So $v$ has $P=\lceil a/2\rceil+\lceil b/2\rceil$ zeros and $w$ has $Q=\lfloor a/2\rfloor+\lfloor b/2\rfloor$, and
$P-Q\in\{0,1,2\}$. Consider the interleaving $u=v_1,w_1,v_2,w_2,\dots,v_m,w_m$; its zeros are $v_1,\dots,v_P$ and $w_1,\dots,w_Q$. If $P=Q$ or $P=Q+1$, the zeros of $u$ are its first
$P+Q$ entries: $u$ is sorted, and the final comparators, which compare $w_i$ with $v_{i+1}$ (neighbours in $u$), change nothing. If $P=Q+2$, $u$ is sorted except that
$w_{Q+1}=1$ comes just before $v_{Q+2}=0$; the comparator on $(w_{Q+1},v_{Q+2})$ swaps them and the others change nothing.
\end{proof}

Sorting by recursive odd--even merges takes the same depth $\frac12p(p+1)$ as the bitonic network but only $\frac n4p(p-1)+n-1$ comparators (Exercise~\ref{ex:oddevencount}): $63$ for
$n=16$, against $80$. The best known network on $16$ wires has $60$ comparators (M.~W.~Green, 1969). And in 1983 Mikl\'os Ajtai, J\'anos Koml\'os and Endre Szemer\'edi proved that
depth $O(\log n)$, and hence size $O(n\log n)$, is possible; their construction uses expander graphs (Tome~C1, Chapter~9), and its constant is so large that Batcher's
$O(\log^2n)$ networks are better for every $n$ that will ever be sorted.

\begin{beauty}
Sorting looks like one problem and is many. Merge sort is optimal up to the round-up of $\log_2k$ at each step; quicksort gambles and almost never loses; selection
needs only a linear number of comparisons, and fewer with a coin than without one; digits can sort faster than comparisons; and a fixed array of comparators,
designed once for all inputs, can be proved correct by looking only at $0$s and $1$s. The zero--one principle is a model of what mathematics does to engineering: it
replaces $n!$ cases by a single idea, here the fact that comparators commute with every monotone function.
\end{beauty}

\section{Exercises}

\begin{exercise}[$\star$]\label{ex:bestmerge}
Show that the \emph{best} case of merge sort on $n=2^p$ elements is $\frac n2p$ comparisons. Which inputs achieve it?
\end{exercise}

\begin{exercise}[$\star\star$]\label{ex:qselectmin}
Quickselect is asked for the minimum ($k=1$). Show that its expected number of comparisons $E(n)$ satisfies $E(n)=n-1+\frac1n\sum_{j=1}^{n-1}E(j)$ and that $E(n)=2n-2H_n$.
\end{exercise}

\begin{exercise}[$\star\star$]\label{ex:groups}
In Theorem~\ref{thm:bfprt}, use groups of $7$ instead of $5$, and show that the time is still linear. With groups of $3$, show that the argument only gives the recurrence
$T(n)\le T(n/3)+T(2n/3)+O(n)$ and hence $O(n\log n)$ by Proposition~1.9(ii).
\end{exercise}

\begin{exercise}[$\star$]\label{ex:radixbase}
Show how to sort $n$ integers in $\{0,\dots,n^3-1\}$ in time $O(n)$ on the word RAM. Why does this not contradict Theorem~1.15?
\end{exercise}

\begin{exercise}[$\star\star$]\label{ex:bucket}
\emph{Bucket sort.} Let $X_1,\dots,X_n$ be independent and uniform on $[0,1)$. Put each $X_i$ into bucket $\lfloor nX_i\rfloor$, sort each bucket by insertion sort and concatenate.
Let $N_j$ be the size of bucket $j$. Show that $\mathbb E\bigl[\sum_jN_j^2\bigr]=2n-1$ and deduce that the expected time is $O(n)$.
\end{exercise}

\begin{exercise}[$\star$]\label{ex:fourwires}
Find a sorting network on $4$ wires with $5$ comparators and depth $3$, prove that it sorts, and prove that $5$ comparators are necessary.
\end{exercise}

\begin{exercise}[$\star\star$]\label{ex:adjacent}
Show that every sorting network on $n$ wires contains, for each $i<n$, a comparator $[i:i+1]$. (Feed it the sorted sequence with the values $i$ and $i+1$ exchanged.)
\end{exercise}

\begin{exercise}[$\star\star$]\label{ex:insertionnet}
The \emph{insertion network} on $n$ wires is built recursively: an insertion network on the first $n-1$ wires, followed by the comparators $[n-1:n],[n-2:n-1],\dots,[1:2]$ in this
order. Prove that it sorts, using the zero--one principle, and show that its comparators can be scheduled in depth $2n-3$.
\end{exercise}

\begin{exercise}[$\star\star$]\label{ex:oddevencount}
Show that the odd--even merge of two sorted lists of length $m=2^q$ uses $M(2^q)=q2^q+1$ comparators, and deduce that odd--even merge sort on $n=2^p$ wires uses $\frac n4p(p-1)+n-1$ comparators.
\end{exercise}

\begin{exercise}[$\star\star$]\label{ex:horses}
Solve the puzzle of the twenty-five horses: seven races suffice to find the three fastest, and at least six races are needed even to find the fastest. (Six races do not suffice for the three
fastest either, but the proof is a longer case analysis.)
\end{exercise}

\section{Hints}

\begin{hint}{ex:bestmerge}
A merge of two lists of length $m$ needs at least $m$ comparisons; when is it exactly $m$?
\end{hint}
\begin{hint}{ex:qselectmin}
The pivot has rank $j+1$ with probability $1/n$, and then the recursion continues on $j$ elements. For the solution, look at $nE(n)-(n-1)E(n-1)$.
\end{hint}
\begin{hint}{ex:groups}
Redo the count of elements below the median of medians with $4$ elements per full group, and redo the induction with the new constants.
\end{hint}
\begin{hint}{ex:radixbase}
Write the numbers in base $n$.
\end{hint}
\begin{hint}{ex:bucket}
$N_j$ is binomial with parameters $n$ and $1/n$; insertion sort on $N$ elements takes $O(N^2+1)$.
\end{hint}
\begin{hint}{ex:fourwires}
Sort two pairs, then compare the two minima and the two maxima, and finish with one comparator in the middle. For the lower bound, use Theorem~1.15.
\end{hint}
\begin{hint}{ex:adjacent}
On the input with $i$ and $i+1$ exchanged, which comparators can change anything?
\end{hint}
\begin{hint}{ex:insertionnet}
The last round of comparators moves the value of wire $n$ up to its place. For the depth, try running comparator $[j:j+1]$ of the stage that inserts wire $k$ at time $2k-j-2$.
\end{hint}
\begin{hint}{ex:oddevencount}
The final round has $m-1$ comparators. Then sum over the merges of odd--even merge sort.
\end{hint}
\begin{hint}{ex:horses}
Race the five groups, then race the five winners. Which horses can still be in the top three? For the lower bound, think of the races as sets of horses: when can one horse be known to beat
all the others?
\end{hint}

\section{Solutions}

\begin{solution}{ex:bestmerge}
\idea Each merge must output the shorter list's elements one comparison at a time; it can do no better than the length of one list.
\writeup Merging lists of lengths $m$ and $m$ makes $2m-r$ comparisons with $r\le m$ (Lemma~\ref{lem:mergecost}), so at least $m$, with equality when all elements of one list are smaller
than all elements of the other. By induction the best case $B(2^p)$ satisfies $B(n)=2B(n/2)+n/2$, so $B(2^p)=\frac n2p$. It is achieved, for instance, by a sorted input: at every merge, the left
list lies entirely below the right one.
\end{solution}

\begin{solution}{ex:qselectmin}
\idea Condition on the rank of the first pivot; the recursion keeps only the elements below it.
\writeup The first round makes $n-1$ comparisons. The pivot has rank $j+1$ with probability $\frac1n$ for $j=0,\dots,n-1$; if $j=0$ it is the minimum, otherwise the search continues among the $j$
smaller elements, with fresh randomness. So $E(n)=n-1+\frac1n\sum_{j=1}^{n-1}E(j)$, $E(1)=0$. Then $nE(n)-(n-1)E(n-1)=n(n-1)-(n-1)(n-2)+E(n-1)=2(n-1)+E(n-1)$, so $n(E(n)-E(n-1))=2(n-1)$, that is
$E(n)-E(n-1)=2-\frac2n$. Summing from $2$ to $n$: $E(n)=2(n-1)-2(H_n-1)=2n-2H_n$.
\end{solution}

\begin{solution}{ex:groups}
\idea The two recursive calls must have total size a fraction less than $1$ of $n$; with groups of $7$ it is $\frac17+\frac57=\frac67$, with groups of $3$ it is exactly $1$.
\writeup With groups of $7$: at least half of the $g=\lceil n/7\rceil$ medians are $\le m$, and each, apart from at most two groups, brings $4$ elements $\le m$; so at least $4(\lceil g/2\rceil-2)\ge\frac{2n}7-8$
elements are $\le m$, and each side has at most $\frac{5n}7+8$. Then $T(n)\le T(\lceil n/7\rceil)+T(\frac{5n}7+8)+an$, and $T(n)\le cn$ holds by induction as soon as
$c(\frac n7+1)+c(\frac{5n}7+8)+an\le cn$, that is $c(\frac n7-9)\ge an$, true for $c\ge14a$ and $n\ge126$. With groups of $3$: each side has at most about $\frac{2n}3$ elements and the sample has $\frac n3$,
so the recurrence is $T(n)\le T(\frac n3)+T(\frac{2n}3)+O(n)$; the two sizes add up to $n$, and Proposition~1.9(ii) gives only $O(n\log n)$. (This is a statement about the argument; Chen and Dumitrescu,
2015, showed that groups of $3$ do work with a cleverer recursion.)
\end{solution}

\begin{solution}{ex:radixbase}
\idea Small keys are short numbers in a large base.
\writeup Write each number in base $n$ with $3$ digits (since it is $<n^3$); a digit is computed with a division, one step on the word RAM since numbers below $n^3$ fit in $O(1)$ words. Radix sort
(Theorem~\ref{thm:radix}) with $B=n$, $d=3$ takes $O(3(n+n))=O(n)$. Theorem~1.15 concerns algorithms that only compare elements; counting sort uses the keys as addresses, which comparisons cannot do.
\end{solution}

\begin{solution}{ex:bucket}
\idea The cost of a bucket is quadratic in its size, so we need the second moment of a binomial variable, not just the mean.
\writeup $N_j$ counts the $i$ with $\lfloor nX_i\rfloor=j$, an event of probability $1/n$, independently for each $i$: $N_j$ is binomial$(n,\frac1n)$, with mean $1$ and variance $1-\frac1n$. So
$\mathbb E[N_j^2]=\operatorname{Var}N_j+(\mathbb EN_j)^2=2-\frac1n$ and $\mathbb E\sum_jN_j^2=n(2-\frac1n)=2n-1$. Insertion sort on $N$ elements takes $O(N^2+1)$ time, so the expected total time is
$O(n)+O(\mathbb E\sum_j(N_j^2+1))=O(n)$.
\end{solution}

\begin{solution}{ex:fourwires}
\idea After sorting two pairs, the overall minimum is one of the two minima and the overall maximum one of the two maxima; only the middle is left.
\writeup Take $[1:2],[3:4]$ (round 1), $[1:3],[2:4]$ (round 2), $[2:3]$ (round 3). After round 1, wires $1,2$ hold $a_1\le a_2$ and wires $3,4$ hold $b_1\le b_2$. The smallest of the four values is
$a_1$ or $b_1$, so $[1:3]$ puts it on wire $1$; the largest is $a_2$ or $b_2$, so $[2:4]$ puts it on wire $4$. Wires $2$ and $3$ then hold the two middle values, in some order, and $[2:3]$ sorts them.
The rounds have depth $3$. A network is a comparison algorithm, so by Theorem~1.15 it needs at least $\lceil\log_224\rceil=5$ comparators.
\end{solution}

\begin{solution}{ex:adjacent}
\idea On an input that is sorted except for one adjacent swap, only a comparator between those two wires can repair it.
\writeup Let the input be $1,2,\dots,n$ with $i$ and $i+1$ exchanged: wire $i$ holds $i+1$ and wire $i+1$ holds $i$. As long as the values are in this configuration, a comparator $[j:j']$ changes
nothing unless its two values are out of order, which happens only for the pair of wires $i,i+1$. So the configuration never changes until a comparator $[i:i+1]$ is applied; since the
network sorts this input, such a comparator exists.
\end{solution}

\begin{solution}{ex:insertionnet}
\idea Like insertion sort: the first $n-1$ values are sorted, and the last one bubbles up to its place.
\writeup By induction and the zero--one principle. After the recursive part, wires $1..n-1$ hold $0^a1^{n-1-a}$ and wire $n$ holds a bit $u$. If $u=1$ the comparators change nothing and the output is
sorted. If $u=0$, the comparator $[n-1:n]$ moves the $0$ to wire $n-1$ (if wire $n-1$ holds a $1$), then $[n-2:n-1]$ moves it up, and so on until it meets the block of $0$s: the output is
$0^{a+1}1^{n-1-a}$. For the depth, number the stages $k=2,\dots,n$ (stage $k$ inserts wire $k$) and run comparator $[j:j+1]$ of stage $k$ at time $2k-j-2$. Within a stage the times are
consecutive in the required order; comparator $[j:j+1]$ of stage $k$ and comparator $[j':j'+1]$ of stage $k+1$ share a wire only if $|j-j'|\le1$, and then the second runs at time
$2k-j'\ge2k-j-1$, later than the first, as needed. Two comparators at the same time $t=2k-j-2=2k'-j'-2$ with $k'=k+1$ have $j'=j+2$, so they touch disjoint wires. The last time is
$2n-1-2=2n-3$ (stage $n$, $j=1$), and the first is $1$ (stage $2$, $j=1$): depth $2n-3$.
\end{solution}

\begin{solution}{ex:oddevencount}
\idea The comparator count of the merge satisfies a merge-sort-like recurrence; then add up the merges level by level.
\writeup The merge of two lists of length $m$ consists of two merges of lists of length $m/2$ and $m-1$ final comparators: $M(m)=2M(m/2)+m-1$, with $M(1)=1$. The formula $M(2^q)=q2^q+1$ holds for $q=0$,
and $2\bigl((q-1)2^{q-1}+1\bigr)+2^q-1=(q-1)2^q+2^q+1=q2^q+1$. Odd--even merge sort on $n=2^p$ wires performs, at level $q=0,\dots,p-1$, $2^{p-1-q}$ merges of two lists of length $2^q$. The total is
\[
\sum_{q=0}^{p-1}2^{p-1-q}\bigl(q2^q+1\bigr)=2^{p-1}\sum_{q=0}^{p-1}q+\sum_{q=0}^{p-1}2^{p-1-q}=2^{p-1}\frac{p(p-1)}2+2^p-1=\frac n4p(p-1)+n-1.
\]
For $n=16$: $48+15=63$.
\end{solution}

\begin{solution}{ex:horses}
\idea After the race of winners, a horse that has three horses known to be faster than it is out; five candidates remain, and one more race settles the top three. For the lower bound, information
only travels between horses that are linked by a chain of races.
\writeup Seven races: race the five groups (races 1--5) and the five group winners (race 6); call these $A_1$ (fastest), $A_2,\dots,A_5$, and let $A_i,B_i,C_i$ be the top three of $A_i$'s group. The fastest
horse is $A_1$. A horse can be second or third only if at most one horse other than $A_1$ is known to be faster: this leaves $A_2,A_3,B_1,C_1,B_2$ (every other horse has at least three known faster
horses, for instance $B_3<A_3<A_2<A_1$ and $C_2<B_2<A_2<A_1$). Race these five (race 7): its top two are second and third overall.

Lower bound. Call two horses linked if they ran in a common race, and consider the connected components of the resulting graph. The results of the races only order horses within a component: if
$X$ and $Y$ are in different components, giving every horse of $Y$'s component a speed larger than all others (keeping their order) is consistent with all results, and so is doing the same for
$X$'s component. So to know the fastest horse, all $25$ horses must be in one component. Each race joins at most $5$ components into one, decreasing their number by at most $4$; starting from $25$
components, reaching $1$ needs at least $24/4=6$ races.
\end{solution}

\section*{Sources and further reading}

\begin{itemize}
\item \emph{The stories.} D.~E.~Knuth, \emph{Von Neumann's first computer program}, Computing Surveys 2 (1970), 247--260; C.~A.~R.~Hoare, \emph{Algorithm 64: Quicksort}, Communications of the ACM 4 (1961), 321,
and \emph{Quicksort}, The Computer Journal 5 (1962), 10--15; Hoare's account of the invention in Moscow is in M.~Freiberger, \emph{Happy birthday Quicksort}, Plus Magazine (University of Cambridge), June 2021, an
interview with Hoare.
\item D.~E.~Knuth, \emph{The Art of Computer Programming}, vol.~3, \emph{Sorting and Searching}, 2nd ed., Addison-Wesley, 1998: Sections~5.2.2 (quicksort), 5.2.4 (merging), 5.2.5 (radix sorting and its history),
5.3.3 (selection) and 5.3.4 (networks, the zero--one principle and Green's $60$-comparator network); D.~E.~Knuth, \emph{Mathematical analysis of algorithms}, Proceedings of IFIP Congress 1971, 19--27, for the average of quickselect.
\item M.~R\'egnier, \emph{A limiting distribution for quicksort}, RAIRO Informatique Th\'eorique 23 (1989), 335--343; U.~R\"osler, \emph{A limit theorem for ``Quicksort''}, RAIRO Informatique Th\'eorique 25 (1991), 85--100.
\item M.~Blum, R.~W.~Floyd, V.~Pratt, R.~L.~Rivest and R.~E.~Tarjan, \emph{Time bounds for selection}, Journal of Computer and System Sciences 7 (1973), 448--461; R.~W.~Floyd and R.~L.~Rivest, \emph{Expected time bounds for
selection}, Communications of the ACM 18 (1975), 165--172; S.~W.~Bent and J.~W.~John, \emph{Finding the median requires $2n$ comparisons}, Proceedings of the 17th ACM STOC (1985), 213--216; R.~Motwani and P.~Raghavan,
\emph{Randomized Algorithms}, Cambridge University Press, 1995, Section~3.3, the source of Theorem~\ref{thm:lazyselect}; K.~Chen and A.~Dumitrescu, \emph{Select with groups of 3 or 4}, Proceedings of WADS 2015, 189--199.
\item K.~E.~Batcher, \emph{Sorting networks and their applications}, Proceedings of the AFIPS Spring Joint Computer Conference 32 (1968), 307--314; M.~Ajtai, J.~Koml\'os and E.~Szemer\'edi, \emph{Sorting in $c\log n$ parallel
steps}, Combinatorica 3 (1983), 1--19; T.~H.~Cormen, C.~E.~Leiserson, R.~L.~Rivest and C.~Stein, \emph{Introduction to Algorithms}, 2nd ed., MIT Press, 2001, Chapter~27, for the half-cleaner presentation.
\end{itemize}
